% main_v4.tex — Vitrine: v4 Current-State Account (consolidated)
% Synthesises the v1 bid aims, the v2 architecture, and the v3 implementation
% plan into a single current-state report, updated with the end-to-end
% validation performed in June 2026. Leaves main.tex (v1), main_v2.tex (v2) and
% main_v3.tex (v3) untouched as historical snapshots.
%
% Self-contained: no biber/glossaries/index, so it compiles cleanly with
% `latexmk -pdf main_v4` or two `pdflatex` passes.
%
% Dr John O'Hare, DreamLab AI Consulting Ltd — June 2026

\documentclass[a4paper,11pt]{report}

\usepackage[utf8]{inputenc}
\usepackage[T1]{fontenc}
\usepackage[british]{babel}
\usepackage{microtype}
\usepackage[a4paper,top=28mm,bottom=28mm,inner=26mm,outer=24mm,headheight=22pt]{geometry}
\usepackage{lmodern}
\usepackage{inconsolata}
\usepackage[dvipsnames,svgnames,x11names]{xcolor}

\definecolor{UoSRed}{HTML}{B71234}
\definecolor{DreamLabAccent}{HTML}{6366F1}
\definecolor{DarkBg}{HTML}{0F172A}
\definecolor{MidGrey}{HTML}{64748B}
\definecolor{LightGrey}{HTML}{F1F5F9}
\definecolor{CodeBg}{HTML}{1E293B}
\definecolor{SuccessGreen}{HTML}{15803D}
\definecolor{WarningAmber}{HTML}{B45309}
\definecolor{DangerRed}{HTML}{B91C1C}

\usepackage{graphicx}
\graphicspath{{figures/}{figures/v2/}{figures/v3/}}
\usepackage{tikz}
\usetikzlibrary{calc,positioning,arrows.meta,shapes.geometric,backgrounds,fit}
\usepackage{booktabs}
\usepackage{tabularx}
\usepackage{longtable}
\usepackage{array}
\newcolumntype{L}[1]{>{\raggedright\arraybackslash}p{#1}}
\newcolumntype{C}[1]{>{\centering\arraybackslash}p{#1}}
\usepackage{amsmath,amssymb}
\usepackage{csquotes}
\usepackage{enumitem}
\setlist{nosep,leftmargin=*}
\usepackage{caption}
\captionsetup{font=small,labelfont={bf,color=UoSRed}}

\usepackage[colorlinks=true,linkcolor=UoSRed,citecolor=DreamLabAccent,urlcolor=DreamLabAccent,
  bookmarks=true,bookmarksnumbered=true,
  pdfauthor={Dr John O'Hare},
  pdftitle={Vitrine: Current-State Account (v4)},
  pdfsubject={Agentic volumetric capture, 3D Gaussian Splatting, USD, Cultural Heritage}]{hyperref}
\usepackage{bookmark}

\usepackage{fancyhdr}
\pagestyle{fancy}
\fancyhf{}
\fancyhead[L]{\small\textcolor{MidGrey}{\leftmark}}
\fancyfoot[R]{\small\textcolor{MidGrey}{\thepage}}
\fancyfoot[C]{\small\textcolor{MidGrey}{Vitrine, DreamLab AI Consulting Ltd}}
\renewcommand{\headrulewidth}{0.4pt}
\renewcommand{\footrulewidth}{0.2pt}
\renewcommand{\headrule}{\hbox to\headwidth{\color{UoSRed}\leaders\hrule height \headrulewidth\hfill}}

\usepackage{listings}
\lstset{basicstyle=\ttfamily\footnotesize,backgroundcolor=\color{LightGrey},frame=single,
  rulecolor=\color{MidGrey},breaklines=true,showstringspaces=false,tabsize=2,
  keywordstyle=\color{DreamLabAccent}\bfseries,commentstyle=\color{MidGrey}\itshape,
  stringstyle=\color{SuccessGreen},extendedchars=true,
  literate={—}{{---}}1 {–}{{--}}1 {→}{{$\rightarrow$}}1 {↔}{{$\leftrightarrow$}}1
    {✓}{{\checkmark}}1 {✗}{{$\times$}}1 {≈}{{$\approx$}}1 {×}{{$\times$}}1
    {‘}{{`}}1 {’}{{'}}1 {“}{{``}}1 {”}{{''}}1}

\usepackage[most]{tcolorbox}
\newtcolorbox{keyfinding}[1][]{enhanced,colback=UoSRed!5,colframe=UoSRed,fonttitle=\bfseries,
  title={Key Finding},left=5mm,boxrule=0.8pt,#1}
\newtcolorbox{recommendation}[1][]{enhanced,colback=DreamLabAccent!6,colframe=DreamLabAccent,
  fonttitle=\bfseries,title={Recommendation},left=5mm,boxrule=0.8pt,#1}
\newtcolorbox{technote}[1][]{enhanced,colback=LightGrey,colframe=MidGrey,fonttitle=\bfseries,
  title={Technical Note},left=5mm,boxrule=0.6pt,#1}

\usepackage{titlesec}
\titleformat{\chapter}[display]{\normalfont\huge\bfseries\color{DarkBg}}
  {\textcolor{UoSRed}{\chaptertitlename\ \thechapter}}{8pt}{\Huge}
\titlespacing*{\chapter}{0pt}{-18pt}{26pt}
\titleformat{\section}{\normalfont\Large\bfseries\color{DarkBg}}{\textcolor{UoSRed}{\thesection}}{1em}{}
\titleformat{\subsection}{\normalfont\large\bfseries\color{DarkBg}}{\textcolor{DreamLabAccent}{\thesubsection}}{1em}{}

% Convenience: status chips
\newcommand{\chip}[2]{\textcolor{#1}{\textbf{#2}}}
\newcommand{\code}[1]{\texttt{#1}}
\newcommand{\validated}{\chip{SuccessGreen}{built \& validated}}
\newcommand{\built}{\chip{SuccessGreen}{built}}
\newcommand{\wip}{\chip{WarningAmber}{in progress}}
\newcommand{\designed}{\chip{MidGrey}{designed}}

\title{Vitrine: Current-State Account}
\author{Dr John O'Hare \\ DreamLab AI Consulting Ltd}
\date{June 2026}

\begin{document}
\sloppy

% ───────────────────────── Title page ─────────────────────────
\begin{titlepage}
\begin{tikzpicture}[remember picture, overlay]
  \fill[DarkBg] (current page.south west) rectangle (current page.north east);
  \foreach \r in {2,3.2,4.4,5.6} {
    \draw[UoSRed, opacity=0.15, line width=0.6pt]
      ($(current page.north west) + (2.5,-3)$) circle (\r);}
  \foreach \i in {0,0.8,...,6.4} {
    \draw[DreamLabAccent, opacity=0.10, line width=0.5pt]
      ($(current page.south east) + (-\i-1, 0.5)$) -- ($(current page.south east) + (-0.5, \i+1)$);}
  \fill[UoSRed] ($(current page.west) + (0,2)$) rectangle ($(current page.east) + (0,2.08)$);
  \fill[DreamLabAccent] ($(current page.west) + (0,1.9)$) rectangle ($(current page.east) + (0,1.94)$);
  \node[anchor=north west, text=UoSRed, font=\footnotesize\sffamily\bfseries, text width=13cm]
    at ($(current page.north west) + (3,-2)$)
    {CURRENT-STATE ACCOUNT \textcolor{MidGrey}{|} END-TO-END VALIDATED};
  \node[anchor=west, text=white, font=\fontsize{40}{44}\selectfont\bfseries\sffamily, text width=15cm]
    at ($(current page.west) + (3,5.6)$) {Vitrine};
  \node[anchor=west, text=white, font=\fontsize{15}{19}\selectfont\sffamily, text width=14cm]
    at ($(current page.west) + (3,2.9)$)
    {An Agentic Pipeline for Volumetric\\ Cultural-Heritage Preservation};
  \node[anchor=west, text=MidGrey, font=\normalsize\sffamily, text width=13cm]
    at ($(current page.west) + (3,-0.2)$)
    {video \textbullet\ COLMAP \textbullet\ 3D Gaussian Splatting \textbullet\ SAM3 \textbullet\ depth-aware object isolation\\
     per-object mesh \textbullet\ native + composed USD scene graph};
  \node[anchor=west, text=white, font=\sffamily, text width=10cm]
    at ($(current page.west) + (3,-2.7)$)
    {\textbf{Dr John O'Hare}\\[2pt] DreamLab AI Consulting Ltd\\[8pt]
     \textcolor{MidGrey}{Pairs with the upstream}\\[2pt] \textbf{LichtFeld Studio} reconstruction engine};
  \node[anchor=west, text=DreamLabAccent, font=\sffamily\bfseries]
    at ($(current page.west) + (3,-5.4)$) {June 2026};
  \node[anchor=south east, text=MidGrey, font=\tiny\sffamily]
    at ($(current page.south east) + (-3,1)$)
    {v4.0 \textcolor{MidGrey}{|} CONSOLIDATED CURRENT-STATE ACCOUNT};
\end{tikzpicture}
\end{titlepage}

% ───────────────────────── Front matter ─────────────────────────
\pagenumbering{roman}

\chapter*{Executive Summary}
\addcontentsline{toc}{chapter}{Executive Summary}
\markboth{Executive Summary}{Executive Summary}

This is the \textbf{single current-state account} of \textbf{Vitrine}, the agentic pipeline that turns a
hand-held walkthrough video of a museum or heritage exhibit into a structured, object-resolved 3D scene.
It consolidates the three prior reports — the original bid (\code{main.tex}), the v2 architecture
(\code{main\_v2.tex}) and the v3 implementation plan (\code{main\_v3.tex}), which remain as historical
snapshots — and updates them with the \textbf{end-to-end validation performed in June 2026}.

\textbf{The headline change since v3:} the pipeline is no longer only \emph{designed} — it has been
\textbf{run end to end on a real reconstructed scene and now produces object-resolved, dual-USD output}.
A single walkthrough now travels the full path: LichtFeld training of a Gaussian-splat field, SAM3
concept segmentation into named objects, a new depth-aware multi-view projection that isolates each
object into its own gaussian subset and ranks them by curatorial significance, per-object meshing, and
\textbf{two USD products} — a native LichtFeld splat USD and a textured composed multi-object scene graph
with preview renders. Eight concrete defects that previously blocked this path were found and fixed, each
surfaced by actually running the pipeline rather than by inspection.

\begin{keyfinding}
Vitrine's core promise --- an \emph{object-resolved} scene, not a monolithic splat --- is now a working
capability, validated on a real 80-frame indoor scene that resolved into five named objects (sculptures,
furniture, walls, floor, ceiling), each isolated into its own gaussian subset and correctly ranked with
the key items (sculptures, furniture) above the structural surfaces.
\end{keyfinding}

\textbf{The honest boundary.} This report uses an explicit \validated{} / \wip{} / \designed{} taxonomy
and does not blur it. Validation was performed on a \emph{reused} 80-frame reconstruction
(\code{milo\_run}), not on a fresh Google-Drive capture driven through ingest and COLMAP; the
fresh-capture path needs Drive credentials and has not been exercised at scale. The per-object meshes are
produced by the always-available \code{gsplat}-to-TSDF backend, \emph{not} by the SOTA single-image hull
generators (TRELLIS.2 / Hunyuan3D-2.1), whose ComfyUI custom-node dependencies are not yet built; and the
generative \emph{recovery} of unseen object faces (FLUX.2) and the local gemma-4 vision agent are staged
but not wired into the loop. Those are reported as \wip{} or \designed{} below, not as done. Quality, too,
is first-pass: SAM3's per-frame detection was sparse and the projection is multi-view-consistent but not
yet depth-tested, so object masks are coherent but over-inclusive.

The report has four chapters: the project and its origin (\S1); the architecture as it now stands (\S2);
the current implementation status, capability by capability, with the June 2026 validation results (\S3);
and what remains, with the decision-record status table and the path to a fresh-capture, SOTA-hull,
fully-agentic run (\S4).

\tableofcontents

% ───────────────────────── Main matter ─────────────────────────
\cleardoublepage
\pagenumbering{arabic}

% ====================================================================
\chapter{The Project and Its Origin}

\section{Why Vitrine exists}
Vitrine began as a University of Salford \emph{AI Knowledge-Exchange Community of Practice} pilot, led from
the MediaCity campus, partnered with DreamLab AI Consulting. The brief — preserved in \code{docs/brief/} —
named the problem plainly: the University has a growing catalogue of immersive, experiential, time-based
arts (Nested Cinema, We Invented the Weekend, Lightwaves) but \emph{no coherent or efficient pipeline for
the preservation and emulation of site-based interactive and immersive experiences}. The pilot's aim was
to test a novel, efficient pipeline that takes existing digital assets — video recordings, still imagery —
and deploys \emph{AI agentic workflows} to produce 3D objects suitable for emulation and preservation
inside a games engine (the Salford MetaCampus, in Unreal).

\section{What the pipeline produces}
That framing maps directly onto what Vitrine builds. From one walkthrough video it produces:
\begin{itemize}
  \item a photoreal \textbf{Gaussian-splat field} of the whole space (the high-fidelity record);
  \item a set of \textbf{named, individually-addressable objects}, each isolated from the field and placed
        at its measured pose (the curatorially-useful decomposition);
  \item a textured \textbf{polygonal environment} mesh and per-object meshes (the games-engine-ready
        geometry); and
  \item a single \textbf{USD scene graph} tying them together with source-video provenance — directly
        importable into Unreal or any USD-aware engine.
\end{itemize}
The work pairs with, and never forks, the upstream \textbf{LichtFeld Studio} reconstruction engine; Vitrine
is the agentic orchestration, segmentation, decomposition and assembly layer on top of it.

% ====================================================================
\chapter{Architecture as It Now Stands}

\section{One walkthrough, one structured scene}
The validated pipeline is a \emph{reconstruct-then-decompose} flow, orchestrated by an in-container agent
over the LichtFeld tool surface. The stages, in order:

\begin{lstlisting}
video -> frames (+ per-image metadata sidecar)
      -> COLMAP SfM  ->  3DGS training (LichtFeld: ImprovedGS+ / MRNF / MCMC)
      -> SAM3 concept segmentation  ->  key-item keyness ranking
      -> depth-aware multi-view mask -> 3D projection (per-object gaussian subsets)
      -> per-object meshing (gsplat-TSDF; CoMe / MILo backends interchangeable)
      -> native USD (LichtFeld convert)  +  composed textured USD (Blender)
\end{lstlisting}

\section{The two-agent separation}
A recurring design principle, unchanged from v3, is the clean separation between the \emph{in-container
orchestrator} (which decides what to run, calling LichtFeld and the pipeline stages as tools) and a
\emph{unified multimodal agent} (a vision-and-reasoning tool it calls for artefact judgement). There is no
hidden state machine: the orchestration is explicit tool-calls over a manifest-driven configuration.

\section{Operational scaffolding}
Underneath sits the v3 operational layer, built and host-validated: a single human-authored
\code{exhibit.toml} manifest with \code{env:} secret indirection and pre-snapshot stripping; a
VRAM-bounded serial model lifecycle; a Docker service mesh (\code{v2g-net}) with service-DNS endpoints; an
agent-controlled ComfyUI client; a state-of-the-art model registry with a host ``idiot-check'' wired into
preflight; and a Rust/Axum web onboarding tool whose secret-containment is verified end-to-end.

% ====================================================================
\chapter{Current Implementation Status}

This chapter is the heart of the report: each capability with its status and, where it was exercised, the
\textbf{June 2026 validation result}. The validation scene throughout is \code{milo\_run}: a real
80-frame indoor capture, reconstructed with COLMAP and trained to a 4M-gaussian splat
(\code{splat\_30000.ply}, LichtFeld \code{igs+} strategy, SH degree 3).

\section{Capability status table}
\begin{center}\small
\begin{tabularx}{\linewidth}{@{}L{0.30\linewidth} c X@{}}
\toprule
\textbf{Capability} & \textbf{Status} & \textbf{Evidence / note (June 2026)} \\
\midrule
LichtFeld training in-container & \validated{} & Runs after the CUDA-13 + vcpkg-OpenUSD + \code{libz-ng} runtime-lib fix; produced \code{splat\_30000.ply}. \\
SAM3 concept segmentation & \validated{} & 5 objects (sculptures, furniture, walls, floor, ceiling) after patching the SAM 3.1 \code{addmm\_act} bf16 dtype regression (sam3 \#507). \\
Depth-aware object isolation (D10) & \validated{} & New multi-view mask$\rightarrow$3D projection: per-object gaussian subsets via COLMAP-camera voting. Isolated sculptures 1{,}080{,}171; furniture 452{,}401; floor 964{,}074; walls 357{,}276; ceiling 11{,}012 gaussians. \\
Key-item keyness ranking (FR-9) & \validated{} & Curatorial priority $\times$ gaussian count ranks sculptures/furniture above structural surfaces despite floor's larger count. \code{min\_object\_gaussians} threshold enforced. \\
Per-object meshing & \validated{} & 5 meshes via \code{gsplat}-to-TSDF after fixing the SH-degree PLY loader (the trained PLY is degree 1$\rightarrow$3; was hardcoded to 45 coefficients). \\
Native USD export & \validated{} & Via the LichtFeld CLI \code{convert} subcommand (headless; no GUI/MCP) $\rightarrow$ \code{scene\_native.usd} (901\,MB splat field). \\
Composed textured USD & \validated{} & Blender assembler $\rightarrow$ \code{scene.usda} + 4 preview renders, after fixing the bake object-selection bug and the Blender 5.0 \code{usd\_export} keyword change. \\
Manifest / lifecycle / service mesh & \validated{} & \code{model\_lifecycle} + \code{comfyui\_control} suites 31/31; manifest loader + idiot-check host-validated (ADR-013). \\
SOTA model registry idiot-check & \validated{} & Wired into preflight; full set (FLUX.2-dev, gemma-4-26B-A4B, TRELLIS.2, Hunyuan3D-2.1, SAM3D) staged on the host. \\
Web onboarding tool & \validated{} & Rust/Axum wizard; secret-containment verified; step-by-step credential instructions (ADR-015). \\
\midrule
SOTA single-image hulls (TRELLIS.2 / Hunyuan3D-2.1) & \wip{} & Weights staged; ComfyUI custom-node dependencies not built. Meshing falls back to gsplat-TSDF. \\
FLUX.2 generative recovery loop & \wip{} & Client built; not yet wired into the per-object hull path (unseen-face inpaint). \\
Local gemma-4 VLM serving & \wip{} & Q8\_0 weights staged; \code{agent-vlm} serving not stood up. \\
\code{v2g:*} USD metadata + per-object composition & \wip{} & Supplied by the custom assembler; native-export \code{customData} parity probe pending. \\
Fresh-capture end-to-end (Drive$\rightarrow$ingest$\rightarrow$COLMAP) & \designed{} & Validated on the reused \code{milo\_run} scene only; Drive ingest needs credentials. \\
Automated acceptance gates + golden fixture & \designed{} & 27-gate harness and golden test scene specified, not yet built. \\
\bottomrule
\end{tabularx}
\end{center}

\section{The validated end-to-end run}
On \code{milo\_run}, a single invocation now traverses segmentation $\rightarrow$ isolation $\rightarrow$
meshing $\rightarrow$ dual-USD assembly and produces: five isolated object PLYs; five per-object meshes; a
901\,MB native splat USD; and a composed, textured \code{scene.usda} with four preview renders. This is
the first time the object-resolved promise has been demonstrated on real data rather than asserted.

\section{Defects found and fixed by running it}
The run was instructive precisely because each stage surfaced a real defect:
\begin{enumerate}
  \item \textbf{SAM3 \#507} --- the SAM 3.1 \code{addmm\_act} fused op casts to bfloat16 and never
        restores, crashing segmentation; patched at the source binding.
  \item \textbf{gsplat SH-degree} --- the PLY loader hardcoded 45 spherical-harmonic coefficients but the
        trained field is SH degree 1 (9); now reads the actual degree and zero-pads.
  \item \textbf{Hunyuan3D kwargs} --- a \code{turbo} keyword crashed the client; kwargs now filtered to
        the constructor signature.
  \item \textbf{LichtFeld runtime} --- the binary needs the host CUDA-13 runtime, vcpkg OpenUSD libs and
        \code{libz-ng}; resolved via \code{LD\_LIBRARY\_PATH} (a Dockerfile fold-in remains).
  \item \textbf{Depth-aware projection (D10)} --- three sub-defects (empty-mask frames inflating the vote,
        absent SH vertex colours, and the world-space alignment) fixed to isolate objects correctly.
  \item \textbf{Native USD path} --- rewired from the never-running MCP server to the headless CLI
        \code{convert}.
  \item \textbf{Blender bake} --- the glTF-import \code{world} root was selected during bake; deselected.
  \item \textbf{Blender USD export} --- Blender 5.0 dropped a \code{usd\_export} keyword; kwargs now
        filtered to the operator's RNA properties.
\end{enumerate}

\section{Known quality limitations}
The isolation mechanism works; its quality is first-pass. SAM3 detected each object in only a minority of
sampled frames, so the multi-view vote runs over few views; the projection is multi-view-consistent (and
parallax filters gross background) but is \emph{not} yet depth-tested per view, so object subsets are
coherent but over-inclusive. Tightening this — more sampled frames, per-view depth occlusion, and denser
SAM3 prompts — is the first quality task, distinct from the functional milestone reported above.

% ====================================================================
\chapter{What Remains}

\section{Decision-record status}
\begin{center}\small
\begin{tabularx}{\linewidth}{@{}c X c@{}}
\toprule
\textbf{ADR} & \textbf{Subject} & \textbf{Status} \\
\midrule
009 & Per-video ingest; per-image metadata sidecar; delete-on-verified-extraction & \wip{} \\
010 & Key-item ranking; depth-aware isolation; hull recovery; pose persistence & \validated{}\,/\,\wip{} \\
011 & USD \code{v2g:*} metadata; native + composed USD; lineage & \validated{}\,/\,\wip{} \\
012 & SOTA tooling: CoMe mesh, ALIKED+LightGlue, TRELLIS.2/Hunyuan, FLUX.2, pins & \wip{} \\
013 & \code{exhibit.toml} manifest; serial lifecycle; \code{v2g-net}; gemma-4 agent & \validated{} \\
014 & Agent-controlled ComfyUI; native USD via CLI convert; recovery loop & \validated{}\,/\,\wip{} \\
015 & Schema-driven web onboarding; hardware-aware selection; Drive write-back & \validated{} \\
\bottomrule
\end{tabularx}
\end{center}
\noindent{\footnotesize\textcolor{MidGrey}{\validated{} = built and exercised on real data \quad \wip{} = core built, integration pending \quad \designed{} = specified, build pending}}

\section{The path to a complete run}
In rough dependency order, the work that turns the validated milestone into a production capability:
\begin{enumerate}
  \item \textbf{Quality:} densify SAM3 sampling and add per-view depth occlusion to the D10 projection.
  \item \textbf{SOTA hulls:} build the TRELLIS.2 / Hunyuan3D-2.1 ComfyUI custom-node dependencies, replacing
        the gsplat-TSDF fallback for key objects.
  \item \textbf{Recovery:} wire the FLUX.2 unseen-face inpaint loop and stand up the local gemma-4 VLM as
        the verify-or-retry judge.
  \item \textbf{Metadata:} populate the \code{v2g:*} per-object USD customData and run the native-export
        parity probe; persist per-object pose.
  \item \textbf{Fresh capture:} exercise the Drive$\rightarrow$ingest$\rightarrow$COLMAP$\rightarrow$\dots
        path end to end on a new capture with credentials.
  \item \textbf{Assurance:} author the golden test fixture and automate the acceptance gates; fold the
        \code{libz-ng} runtime lib and host-staged weights into the container image.
\end{enumerate}

\begin{recommendation}
Vitrine has crossed the line from \emph{designed} to \emph{demonstrated}: a real capture now becomes an
object-resolved, dual-USD scene. The next increment of value is \emph{quality and SOTA hulls} on the same
validated path, followed by a fresh-capture run --- not new architecture. The honest gaps above are the
backlog, in priority order.
\end{recommendation}

\end{document}
